\chapter{Literature Review}\label{ch:2}

\section{Ejection fraction and its ECG signature}\label{sec:2.1}

Clinical guidelines divide LVEF into normal, mildly reduced and reduced ranges \cite{ref1,ref2}, and echocardiography is the standard way to measure it \cite{ref5}. The measurement itself is not exact: two echocardiograms of the same patient may disagree by around 10 EF points \cite{ref6}, which is close to the width of a single class in this project (10--13 points). Several ECG findings are known to accompany a weak left ventricle: a prolonged QRS or bundle-branch block \cite{ref7}, a fragmented QRS \cite{ref8}, a wide angle between the QRS and T vectors \cite{ref9} and uneven repolarisation \cite{ref10,ref11}.

\section{AI on the ECG for LV dysfunction}\label{sec:2.2}

On digital 12-lead ECGs, a convolutional network flagged EF $\le35$ \% with an AUC of 0.93 \cite{ref3}, and in a randomised primary-care trial it led to more new diagnoses \cite{ref4}. Fewer studies grade the EF instead of flagging it: Vaid et al. predicted three EF bands and also regressed the EF value, using 715,890 ECGs \cite{ref12}. Working from ECG \emph{images} rather than signals, a model trained on 385,601 ECGs reached AUROC 0.90--0.95 for EF $<40$ \% \cite{ref13}, and turning ECG images back into signals is itself an active research topic \cite{ref14,ref15,ref16}. Models that predict several structural heart diseases together improved as more input types were added \cite{ref17}, and single-lead and wearable devices have carried the screening idea outside the hospital \cite{ref18,ref19,ref20}.

\section{Foundation models and public data}\label{sec:2.3}

{\small
\begin{longtable}{>{\raggedright\arraybackslash\hspace{0pt}}p{\dimexpr 0.165\linewidth-2\tabcolsep\relax}>{\raggedright\arraybackslash\hspace{0pt}}p{\dimexpr 0.239\linewidth-2\tabcolsep\relax}>{\raggedright\arraybackslash\hspace{0pt}}p{\dimexpr 0.150\linewidth-2\tabcolsep\relax}>{\raggedright\arraybackslash\hspace{0pt}}p{\dimexpr 0.177\linewidth-2\tabcolsep\relax}>{\raggedright\arraybackslash\hspace{0pt}}p{\dimexpr 0.267\linewidth-2\tabcolsep\relax}}
\caption{Relevant studies: data, input, task and main reported result (as printed in each abstract).}\label{tab:2.1}\\
\toprule
\textbf{Study} & \textbf{Data} & \textbf{Input} & \textbf{Task} & \textbf{Main reported result} \\
\midrule
\endfirsthead
\multicolumn{5}{l}{\small\emph{Table~\thetable\ (continued)}}\\
\toprule
\textbf{Study} & \textbf{Data} & \textbf{Input} & \textbf{Task} & \textbf{Main reported result} \\
\midrule
\endhead
\bottomrule
\endlastfoot
Attia et al. \cite{ref3} & 44,959 train / 52,870 test patients, Mayo Clinic & digital 12-lead & EF $\le35$ \% & AUC 0.93, accuracy 85.7 \% \\
Attia et al. \cite{ref32} & 3,874 prospective patients & digital 12-lead & EF $\le35$ \% & AUC 0.918; 39.8 \% of false positives had EF 36--50 \% \\
Yao et al. \cite{ref4} & 22,641 adults, 120 care teams (RCT) & digital 12-lead & new low-EF diagnosis & diagnoses 1.6 \% $\rightarrow$ 2.1 \% (OR 1.32) \\
Adedinsewo et al. \cite{ref33} & 1,606 emergency patients & digital 12-lead & EF $\le35$ \% / $<50$ \% & AUROC 0.89 / 0.85 \\
Vaid et al. \cite{ref12} & 715,890 ECGs, 147,636 patients, 5 hospitals & digital 12-lead & EF 3 classes + regression & AUC 0.94 / 0.82 / 0.89 (internal); MAE 5.8--6.1 \\
Kwon et al. \cite{ref34} & 55,163 ECGs, 2 hospitals & ECG features + demographics & HFrEF (EF $\le40$ \%) & AUROC 0.843 internal, 0.889 external \\
Sangha et al. \cite{ref35} & 2.2 million ECGs plotted as images & ECG images & 6 rhythm / conduction labels & mean AUROC 0.99 (external 0.97) \\
Sangha et al. \cite{ref13} & 385,601 ECGs, Yale + external sites & ECG images & EF $<40$ \% & AUROC 0.91 internal, 0.90--0.95 external \\
Ulloa-Cerna et al. \cite{ref17} & 2.2 million ECGs, 484,765 adults & ECG + age, sex & composite structural disease & AUROC 0.91 \\
Elias et al. \cite{ref36} & 77,163 patients & digital 12-lead & moderate / severe AS, AR, MR & AUROC 0.88 / 0.77 / 0.83 \\
Poterucha et al. \cite{ref23} & $>1$ million ECG-- imaging pairs & digital 12-lead & structural heart disease & outperformed cardiologists; data and weights released \\
Hughes et al. \cite{ref24} & 100,000 ECGs, 36,286 patients (public) & digital, 250 Hz & composite + 11 labels & composite AUROC 82.0 \\
McKeen et al. \cite{ref22} & 1.5 million ECGs (pre-training) & digital 12-lead & EF $\le40$ \% and others & AUROC 0.929 for EF $\le40$ \% \\
Li et al. \cite{ref21} & 10.8 million ECGs, 150 labels & digital 12-lead & broad diagnosis & AUROC $>0.95$ for 80 diagnoses \\
\midrule
\textbf{This work} & 837 development printouts, one centre; EchoNext 4,827 & printout images (and derived traces) & \textbf{four-class EF} & 0.558 all patients; 0.696 (48 \% answered); heart-failure screen AUC up to 0.835 \\
\end{longtable}
}


Foundation models change what a small dataset can achieve. ECGFounder learnt from 10.8 million ECGs and adapts to new tasks with few labelled examples \cite{ref21}, and ECG-FM reported AUROC 0.929 for EF $\le40$ \% \cite{ref22}. The EchoNext project made paired ECG--echocardiography data and trained models public, and its models outperformed cardiologists \cite{ref23,ref24}. Earlier one-dimensional networks had already matched cardiologists in rhythm classification \cite{ref25} and automated 12-lead interpretation \cite{ref26}. For images, this project uses the ResNet \cite{ref27}, EfficientNet \cite{ref28}, MobileNetV3 \cite{ref29} and ConvNeXt \cite{ref30} backbones; for sequences, the LSTM \cite{ref31} and convolutional LSTM variants.

\section{Research gaps addressed in Phase I}\label{sec:2.4}

\begin{itemize}
  \item \textbf{Graded EF from printouts.} Most published models answer a yes/no question on digital signals. Four-class grading that includes Mild, from printout images, had not been reported, and Mild in particular needs more data than single-centre studies usually have.
  \item \textbf{Small single-centre data.} Methods for a few hundred labelled printouts, evaluated with nested cross-validation, are rarely studied.
  \item \textbf{Transfer to the printout format.} Public ECG--echo data had not been redrawn in a hospital's printout format for pre-training.
  \item \textbf{Confident subsets.} Triage that answers only confident cases and refers the rest is seldom evaluated for graded EF.
\end{itemize}
